%% P24 --- Sumy zmiennych losowych: centralne twierdzenie graniczne i metoda Monte Carlo
%%
%% Zalążek. Poniższy opis jest kontraktem tego programu: co ma uzasadnić i co
%% ma dać czytelnikowi. Napisanie programu oznacza usunięcie bloku
%% \programstub{} i zastąpienie go programem. Jeśli po przerobieniu materiału
%% opis okaże się błędny, popraw go i odnotuj sprzeczność w CLAUDE.md w sekcji
%% *Resolved questions*.

\program{Sumy zmiennych losowych: centralne twierdzenie graniczne i metoda Monte Carlo}
\label{prog:P24}

\programstub{%
\textit{[Opis do przetłumaczenia --- patrz wydanie angielskie.]}

Argument: variances of independent quantities add; averages concentrate; the
error of an average falls as \code{1/sqrt(n)} and that single rate governs
an astonishing amount of practice. Payoff: the derivation of the
\code{1/sqrt(d\_k)} in attention. A dot product of two
\code{d\_k}-dimensional vectors with independent unit-variance entries has
variance \code{d\_k}, so its standard deviation grows as \code{sqrt(d\_k)};
feed that into a softmax and it saturates, and the gradient dies; divide by
\code{sqrt(d\_k)} and the logits have unit variance regardless of head size.
The scaling is a variance correction and nothing else, and the reader
derives it here and measures it in P31. Also: Monte Carlo error is the same
\code{1/sqrt(n)}, so the number of samples needed to halve an error bar is
four times as many --- the fact that prices every evaluation run in the
book.

\medskip
\textbf{Planowana długość:} 55 ramek.
}
